\begin{textsample}{POS Dim 3 – df – Score 11.00 – df/df\_inf\_11.txt}  \label{ex:f3_pos_107}
7.4. \textbf{Rotina} É importante enfatizar que a \textbf{rotina} é apenas um dos elementos que compõem o cotidiano.
Geralmente, a \textbf{rotina} abrange recepção, roda de conversa, calendário, clima, alimentação, \textbf{higiene}, atividades de \textbf{pintura} e desenho, \textbf{descanso}, \textbf{brincadeira} livre ou dirigida, narração de histórias, entre outras ações.
Ao planejar a \textbf{rotina} da turma, o professor deve considerar os elementos : materiais, espaços e tempos, bem como os sujeitos que estarão envolvidos nas atividades, pois tudo deve adequar -se à realidade das \textbf{crianças}.
A \textbf{rotina} pode ser o caminho para evitar atividades esvaziadas de sentido, rituais repetitivos, reprodução de regras e fazeres automatizados.
Para tanto, é fundamental que a \textbf{rotina} seja dinâmica e flexível.
Barbosa ( 2006 ) aponta que a \textbf{rotina} inflexível e desinteressante pode vir a ser “ uma tecnologia de alienação ” se não forem levados em consideração o ritmo, a participação, a relação com o mundo, a realização, a fruição, a liberdade, a consciência, a imaginação e as diversas formas de sociabilidade dos sujeitos nela envolvidos.
A \textbf{rotina} é uma forma de organizar o coletivo \textbf{infantil} \textbf{diário} e, concomitantemente, espelha a Proposta Pedagógica da instituição de Educação \textbf{Infantil}.
Ela é capaz ainda de apresentar quais as concepções de educação, de \textbf{criança} e de infância que se materializam no cotidiano educativo.
Com o estabelecimento de objetivos claros e coerentes, a \textbf{rotina} promove aprendizagens, desenvolve a autonomia e a identidade, propicia o movimento corporal, a estimulação dos sentidos, a sensação de segurança e confiança, o suprimento das necessidades biológicas ( alimentação, \textbf{higiene} e repouso ), isso porque contém elementos que devem proporcionar o bem-estar e o desenvolvimento integral da \textbf{criança}.
No caso da jornada em tempo integral, sugere -se que, no período da manhã, sejam incluídas atividades físicas, observando o tempo e a intensidade de calor ou frio.
Já no período da tarde, podem ocorrer atividades como sono ou repouso e banho, ou seja, práticas sociais que envolvem as necessidades vitais dos seres humanos.
Nas jornadas de tempo parcial, por serem mais curtas, tais práticas sociais aparecem com \textbf{menor} frequência, ainda que também estejam presentes.
É essencial abrir espaço e reservar tempo para as \textbf{brincadeiras}, sejam livres ou dirigidas, isso em contextos de Educação \textbf{Infantil} de tempo integral ou parcial.
Vale destacar que as ações da \textbf{rotina} devem se pautar nas necessidades das \textbf{crianças}, e não nas relações de trabalho dos \textbf{adultos}.
Os horários de lanche, almoço, \textbf{limpeza} das salas, funcionamento da cozinha, ou seja, as atividades relacionadas às \textbf{crianças} precisam estar sintonizadas com suas próprias necessidades.
Por vezes, as \textbf{crianças} querem ou propõem outros elementos que transgridam as formalidades da \textbf{rotina}, das jornadas integrais ou parciais, dos momentos instituídos pelos profissionais da educação, sejam no sono, na alimentação, na \textbf{higiene}, na “ hora da atividade ”, nas \textbf{brincadeiras}, entre outros.
A partir da observação, é possível detectar como as \textbf{crianças} vivem o cotidiano da instituição de Educação \textbf{Infantil}.
Esses sinais das \textbf{crianças} ajudam a apontar possibilidades que não se limitam às \textbf{rotinas} formalizadas e ainda oferecem subsídios para trazer à tona a valorização da infância em suas relações e práticas.
Cresce a relevância de um planejamento cuidadoso, flexível, reflexivo que minimize o perigo da \textbf{rotina} ser monótona, distante e vazia de sentido para as \textbf{crianças} e até para os profissionais da educação.

% matched lemmas: adulto, brincadeira, criança, descanso, diário, higiene, infantil, limpeza, pequeno, pintura, rotina
\end{textsample}
